\begin{declarations}

\begin{contributions}
\ifidentificado
DJC: concepção do estudo, metodologia, curadoria dos dados, \textit{software} (\textit{scripts} de análise), análise formal, visualização e redação do manuscrito (rascunho original, revisão e edição). ODRF e DSP: supervisão, revisão crítica do manuscrito (revisão e edição) e aprovação da versão final. Todas as pessoas autoras leram e aprovaram o manuscrito final.
\else
Omitido para revisão.
\fi
\end{contributions}

\begin{interests}
As pessoas autoras declaram que não têm nenhum conflito de interesses.
\end{interests}

\begin{materials}
Os dados analisados são públicos e foram obtidos do INEP (IDEB 2025, microdados do Censo Escolar 2025 e Sinopse Estatística da Educação Básica 2025) e do IBGE (Censo Demográfico 2022, Sidra, tabelas 10295 e 10296). Os \textit{scripts} de ingestão, transformação e análise, com instruções de execução e o ambiente Python reprodutível, estão disponíveis em
\ifidentificado
\url{https://github.com/darlondjc/desempenho-educacional-maranhao}.
\else
repositório anonimizado para a revisão: \url{https://anonymous.4open.science/r/desempenho-educacional-maranhao-97FB/}.
\fi
\end{materials}

\begin{aiusagedeclaration}
Foram utilizadas ferramentas de inteligência artificial generativa nas seguintes etapas: o Claude (Anthropic), na correção e revisão dos \textit{scripts} Python da análise, com cada alteração validada pela autoria; o Claude e o ChatGPT (OpenAI), na revisão da redação e da clareza do texto; o Claude e o SciSpace AI Research (PubGenius Inc.), na busca de trabalhos relevantes e na verificação de metadados de citações, com cada referência conferida na fonte original; e o Claude, na adaptação do manuscrito ao modelo da RBIE, o que incluiu a reorganização das seções, a tradução dos resumos e a redação de versões preliminares de trechos de ligação a partir do conteúdo produzido pela autoria, todos revisados pelas pessoas autoras. Os dados, a análise e a interpretação dos resultados não foram gerados por essas ferramentas. Elas não figuram como autoras, e as pessoas autoras assumem integral responsabilidade por todo o conteúdo do manuscrito.
\end{aiusagedeclaration}

\begin{furtherinformation}
O estudo usa apenas dados secundários públicos, agregados por município, e dispensa apreciação por Comitê de Ética em Pesquisa (ver Seção~\ref{sec:metodos}).
\end{furtherinformation}

\end{declarations}
